%==============================================================================
%  RiskGuard Pro - Presentation de fin d'annee (Beamer)
%  Theme: bleu marine / blanc / accents dores
%  Compilation : pdflatex riskguard_pro.tex  (x2 pour les references)
%  Paquets requis (tous dans TeX Live): beamer, tikz, pgfplots, xcolor,
%                                        booktabs, fontawesome5, qrcode
%
%  Police "moderne" (Poppins/Montserrat) : decommenter le bloc fontspec
%  et compiler avec lualatex ou xelatex au lieu de pdflatex.
%==============================================================================
\documentclass[aspectratio=169,10pt]{beamer}

\usepackage[utf8]{inputenc}
\usepackage[T1]{fontenc}
\usepackage[french]{babel}
\usepackage{helvet}
\renewcommand{\familydefault}{\sfdefault}

% --- Pour une vraie police Poppins/Montserrat (XeLaTeX/LuaLaTeX) ---
% \usepackage{fontspec}
% \setsansfont{Montserrat}
% \renewcommand{\familydefault}{\sfdefault}

\usepackage{tikz}
\usepackage{pgfplots}
\pgfplotsset{compat=1.18}
\usetikzlibrary{positioning,shapes.geometric,arrows.meta,calc,matrix,shadings,backgrounds,fit}
\usepackage{booktabs}
\usepackage{fontawesome5}
\usepackage{qrcode}

%------------------------------------------------------------------------------
% PALETTE
%------------------------------------------------------------------------------
\definecolor{navy}{HTML}{1A237E}
\definecolor{navydark}{HTML}{0A1238}
\definecolor{navymid}{HTML}{0F1A4A}
\definecolor{card}{HTML}{18224F}
\definecolor{gold}{HTML}{FFD700}
\definecolor{goldsoft}{HTML}{FFE566}
\definecolor{golddeep}{HTML}{E6B800}
\definecolor{ink}{HTML}{EAF0FF}
\definecolor{muted}{HTML}{AAB4E0}
\definecolor{rgreen}{HTML}{2ECC71}
\definecolor{ryellow}{HTML}{F1C40F}
\definecolor{rorange}{HTML}{E67E22}
\definecolor{rred}{HTML}{E74C3C}
\definecolor{rblue}{HTML}{3949AB}
\definecolor{rbluel}{HTML}{5C6BC0}

%------------------------------------------------------------------------------
% THEME
%------------------------------------------------------------------------------
\usetheme{default}
\usefonttheme{professionalfonts}
\setbeamertemplate{navigation symbols}{}

\setbeamercolor{background canvas}{bg=navydark}
\setbeamercolor{normal text}{fg=ink}
\setbeamercolor{frametitle}{fg=white}
\setbeamercolor{itemize item}{fg=gold}
\setbeamercolor{itemize subitem}{fg=goldsoft}
\setbeamercolor{enumerate item}{fg=gold}
\setbeamercolor{block title}{fg=navydark,bg=gold}
\setbeamercolor{block body}{fg=ink,bg=card}

\setbeamerfont{frametitle}{size=\Large,series=\bfseries}
\setbeamerfont{title}{size=\Huge,series=\bfseries}
\setbeamertemplate{itemize item}{\small\faCircle}
\setbeamertemplate{itemize subitem}{\scriptsize\faAngleRight}

% --- Titre de frame avec liseré doré ---
\setbeamertemplate{frametitle}{%
  \vspace{0.6em}%
  {\color{goldsoft}\scriptsize\bfseries\MakeUppercase{\insertshortauthor}}\par
  \vspace{0.15em}%
  {\color{white}\Large\bfseries\insertframetitle}\par
  \vspace{0.25em}%
  {\color{gold}\rule{2.2em}{2.4pt}}\par
  \vspace{0.4em}%
}

% --- Bas de page ---
\setbeamertemplate{footline}{%
  \hbox{%
    \begin{beamercolorbox}[wd=\paperwidth,ht=2.4ex,dp=1.2ex,leftskip=1.2em,rightskip=1.2em]{}%
      {\color{gold}\faShieldAlt}\ {\color{ink}\footnotesize\bfseries RiskGuard Pro}%
      \hfill
      {\color{muted}\footnotesize\insertframenumber\,/\,\inserttotalframenumber}%
    \end{beamercolorbox}%
  }%
  \vspace{1pt}\hbox{\color{gold}\rule{\paperwidth}{1.5pt}}%
}

% Pas de "short author" affiche en haut sauf via kicker
\renewcommand{\insertshortauthor}{}

%------------------------------------------------------------------------------
% COMMANDES UTILITAIRES
%------------------------------------------------------------------------------
% kicker = sur-titre dore (passe via \shortauthor du frametitle)
\newcommand{\kick}[1]{\renewcommand{\insertshortauthor}{#1}}

% Badge de timing en haut a droite
\newcommand{\timing}[1]{%
  \begin{tikzpicture}[remember picture,overlay]
    \node[anchor=north east,xshift=-1.1em,yshift=-1.0em,
          draw=golddeep,rounded corners=8pt,inner sep=4pt,
          text=gold,font=\scriptsize] at (current page.north east) {#1};
  \end{tikzpicture}}

% Carte (boite arrondie)
\newcommand{\rgcard}[2]{%
  \begin{tikzpicture}
    \node[fill=card,rounded corners=8pt,draw=navy,line width=0.6pt,
          inner sep=10pt,text=ink,text width=#1] {#2};
  \end{tikzpicture}}

% Tuile "feature"
\newcommand{\feat}[3]{%
  \begin{tikzpicture}
    \node[fill=card,rounded corners=8pt,inner sep=8pt,text width=0.93\linewidth] {%
      {\color{gold}#1}\quad{\color{white}\bfseries #2}\\[2pt]
      {\color{muted}\small #3}};
  \end{tikzpicture}}

% Style global des graphiques pgfplots (fond sombre)
\pgfplotsset{
  rgaxis/.style={
    width=\linewidth,
    axis line style={muted},
    tick label style={muted,font=\scriptsize},
    label style={muted,font=\scriptsize},
    every axis title/.style={muted,font=\small},
    grid style={navy!60},
    tick style={muted},
  }
}

%------------------------------------------------------------------------------
\title{RiskGuard Pro}
\author{}
\date{2025--2026}
%------------------------------------------------------------------------------
\begin{document}

%=========================== SLIDE 1 - GARDE ==================================
{
\setbeamertemplate{footline}{}
\begin{frame}[plain]
\begin{tikzpicture}[remember picture,overlay]
  \fill[navydark] (current page.south west) rectangle (current page.north east);
  % bouclier
  \begin{scope}[shift={($(current page.center)+(0,2.0)$)}]
    \fill[gold] (0,1.1) -- (1.0,0.75) -- (1.0,-0.1) -- (0,-1.05) -- (-1.0,-0.1) -- (-1.0,0.75) -- cycle;
    \draw[navydark,line width=3pt,line cap=round,line join=round] (-0.45,0.0) -- (-0.1,-0.35) -- (0.55,0.45);
  \end{scope}
\end{tikzpicture}

\centering
\vspace{3.6cm}
{\color{gold}\scriptsize\bfseries PROJET DE FIN D'ANNÉE\quad\textbullet\quad 2025--2026}\\[0.9em]
{\color{white}\fontsize{42}{46}\selectfont\bfseries Risk\textcolor{gold}{Guard} Pro}\\[0.3em]
{\color{gold}\rule{4cm}{2.5pt}}\\[0.8em]
{\color{muted}\large Plateforme Intelligente de Détection de Risque Bancaire pour Entreprises}\\[1.6em]

\begin{tabular}{ccc}
  {\color{white}\bfseries Étudiant(s)} & {\color{white}\bfseries Encadrant} & {\color{white}\bfseries Université}\\[2pt]
  {\color{goldsoft}\itshape [ à compléter ]} & {\color{goldsoft}\itshape [ à compléter ]} & {\color{goldsoft}\itshape [ logo / nom ]}\\
\end{tabular}\\[1.4em]

\tikz\node[draw=gold,rounded corners=14pt,inner sep=6pt,text=gold,font=\small\bfseries]
  {Année universitaire 2025 -- 2026};
\end{frame}
}

%=========================== SLIDE 2 - SOMMAIRE ===============================
\kick{Plan de la présentation}
\begin{frame}{Sommaire}
\timing{30 s}
\vspace{0.3em}
\begin{columns}[T]
\column{0.5\textwidth}
  \begin{enumerate}
    \item Contexte \& Problématique
    \item Objectifs du projet
    \item Données \& Prétraitement
    \item Données synthétiques (SMOTE + Bruit Gaussien)
    \item Comparaison des modèles ML
  \end{enumerate}
\column{0.5\textwidth}
  \begin{enumerate}
    \setcounter{enumi}{5}
    \item Pourquoi XGBoost ?
    \item Architecture technique
    \item Démonstration
    \item Résultats \& Métriques
    \item Conclusion \& Perspectives
  \end{enumerate}
\end{columns}
\end{frame}

%=========================== SLIDE 3 - CONTEXTE ===============================
\kick{Pourquoi ce projet ?}
\begin{frame}{Contexte \& Problématique}
\timing{1 min}
\begin{columns}[T]
\column{0.56\textwidth}
  \begin{itemize}
    \item Le secteur bancaire marocain fait face à un \textbf{risque croissant de défaillance} d'entreprises.
    \item Les méthodes traditionnelles (scoring manuel, ratios simples) sont \textbf{limitées et subjectives}.
    \item Besoin d'un outil \textbf{automatisé, fiable et rapide} pour classifier le niveau de risque.
  \end{itemize}
  \vspace{0.5em}
  \begin{block}{\faQuestionCircle\ Question centrale}
    Comment prédire la \textbf{probabilité de faillite} d'une entreprise à partir de ses \textbf{ratios financiers} ?
  \end{block}
\column{0.44\textwidth}
  \centering
  {\color{gold}\small Hausse des défaillances (indice)}\\[2pt]
  \begin{tikzpicture}
    \begin{axis}[rgaxis,height=4.4cm,ymin=0,ymax=90,
      symbolic x coords={2018,2019,2020,2021,2022,2023,2024,2025},
      xtick={2018,2021,2024}, ytick=\empty,ymajorgrids]
      \addplot[gold!50,fill=gold!12,draw=none] coordinates
        {(2018,22)(2019,28)(2020,26)(2021,35)(2022,41)(2023,52)(2024,63)(2025,78)}
        \closedcycle;
      \addplot[gold,line width=1.5pt,mark=*,mark options={fill=navydark}]
        coordinates {(2018,22)(2019,28)(2020,26)(2021,35)(2022,41)(2023,52)(2024,63)(2025,78)};
    \end{axis}
  \end{tikzpicture}
\end{columns}
\end{frame}

%=========================== SLIDE 4 - OBJECTIFS ==============================
\kick{Ce que nous voulons livrer}
\begin{frame}{Objectifs du projet}
\timing{1 min}
\vspace{0.4em}
\begin{columns}[T]
\column{0.5\textwidth}
  \feat{\faProjectDiagram}{Pipeline ML complet}{De la donnée brute à la prédiction temps réel.}\\[6pt]
  \feat{\faLayerGroup}{4 niveaux de risque}{Faible -- Modéré -- Élevé -- Critique.}\\[6pt]
  \feat{\faDesktop}{Interface web intuitive}{Pensée pour les analystes bancaires.}
\column{0.5\textwidth}
  \feat{\faSearch}{Interprétabilité (SHAP)}{Expliquer chaque décision du modèle.}\\[6pt]
  \feat{\faFileCsv}{Analyse unitaire \& batch}{Saisie manuelle ou import CSV/Excel.}\\[6pt]
  \feat{\faBolt}{Réponse temps réel}{Prédiction en moins de 100 ms.}
\end{columns}
\end{frame}

%=========================== SLIDE 5 - DATASET ================================
\kick{Le jeu de données}
\begin{frame}{Données utilisées}
\timing{1 min}
\begin{columns}[T]
\column{0.56\textwidth}
  \begin{itemize}
    \item \textbf{Source :} « Company Bankruptcy Prediction » (Taiwan Economic Journal).
    \item \textbf{96 colonnes} de ratios financiers : ROA, ROE, liquidité, endettement, trésorerie\dots
    \item \textbf{Cible :} \texttt{Bankrupt?} -- 0 = sain, 1 = défaillant.
    \item Environ \textbf{6 800 observations} initiales.
  \end{itemize}
  \vspace{0.3em}
  \begin{tikzpicture}
    \node[fill=rred!18,draw=rred!70,rounded corners=8pt,inner sep=9pt,
          text width=0.93\linewidth,text=ink]{%
      {\color{rred}\faExclamationTriangle\ \bfseries Problème majeur}\\[3pt]
      \textbf{Déséquilibre extrême} : \textasciitilde3\% de faillites contre 97\% d'entreprises saines.};
  \end{tikzpicture}
\column{0.44\textwidth}
  \centering
  {\color{gold}\small Déséquilibre des classes}\\[2pt]
  \begin{tikzpicture}
    \begin{axis}[rgaxis,height=4.6cm,ybar,ymin=0,ymax=100,bar width=34pt,
      symbolic x coords={Sains,Faillite},xtick=data,
      nodes near coords,nodes near coords style={text=white,font=\small\bfseries},
      ytick=\empty]
      \addplot[fill=rblue,draw=none] coordinates {(Sains,97)};
      \addplot[fill=rred,draw=none] coordinates {(Faillite,3)};
    \end{axis}
  \end{tikzpicture}
\end{columns}
\end{frame}

%=========================== SLIDE 6 - DATA CLEANING ==========================
\kick{Prétraitement}
\begin{frame}{Data Cleaning}
\timing{1 min 30}
\begin{columns}[T]
\column{0.58\textwidth}
  \small
  \begin{enumerate}
    \item Détection \& suppression des \textbf{doublons}.
    \item Valeurs manquantes : \textbf{imputation par la médiane}.
    \item Outliers via \textbf{IQR} -- capping percentiles 1\% / 99\%.
    \item Suppression des features à \textbf{variance quasi-nulle}.
    \item Corrélation : suppression si \textbf{corr. $>$ 0.95}.
    \item \textbf{Normalisation} des noms de colonnes.
  \end{enumerate}
\column{0.42\textwidth}
  \centering
  {\color{gold}\small Réduction des features}\\[2pt]
  \begin{tikzpicture}
    \begin{axis}[rgaxis,height=3.6cm,ybar,ymin=0,ymax=110,bar width=30pt,
      symbolic x coords={Avant,Après},xtick=data,
      nodes near coords,nodes near coords style={text=white,font=\small\bfseries},
      ytick=\empty]
      \addplot[fill=rbluel,draw=none] coordinates {(Avant,96)};
      \addplot[fill=gold,draw=none] coordinates {(Après,65)};
    \end{axis}
  \end{tikzpicture}
  {\color{goldsoft}\footnotesize 96 features $\rightarrow$ \textbf{\textasciitilde65} pertinentes}
\end{columns}
\end{frame}

%=========================== SLIDE 7 - POURQUOI SYNTH =========================
\kick{Le piège du déséquilibre}
\begin{frame}{Pourquoi générer des données synthétiques ?}
\timing{1 min}
\begin{columns}[T]
\column{0.52\textwidth}
  \begin{itemize}
    \item Avec \textasciitilde3\% de positifs, le modèle apprend à \textbf{toujours prédire « sain »}.
    \item \textbf{Accuracy trompeuse} : 97\% en prédisant toujours 0 $\rightarrow$ inutile !
    \item Le \textbf{recall} sur la classe « faillite » est quasi nul.
    \item \textbf{Solution :} augmenter la classe minoritaire avec des données synthétiques réalistes.
  \end{itemize}
\column{0.48\textwidth}
  \centering {\color{gold}\small Matrice de confusion}\\[4pt]
  \begin{columns}[T]
  \column{0.5\textwidth}\centering
    {\color{muted}\scriptsize Sans rééquilibrage}\\[2pt]
    \begin{tikzpicture}[scale=0.62,every node/.style={font=\footnotesize\bfseries,text=white}]
      \fill[rgreen!70] (0,1) rectangle (1,2); \node at (0.5,1.5){97};
      \fill[rred!25]   (1,1) rectangle (2,2); \node at (1.5,1.5){0};
      \fill[rred!45]   (0,0) rectangle (1,1); \node at (0.5,0.5){3};
      \fill[rred!25]   (1,0) rectangle (2,1); \node at (1.5,0.5){0};
    \end{tikzpicture}
  \column{0.5\textwidth}\centering
    {\color{muted}\scriptsize Après rééquilibrage}\\[2pt]
    \begin{tikzpicture}[scale=0.62,every node/.style={font=\footnotesize\bfseries,text=white}]
      \fill[rgreen!75] (0,1) rectangle (1,2); \node at (0.5,1.5){88};
      \fill[rred!35]   (1,1) rectangle (2,2); \node at (1.5,1.5){9};
      \fill[rred!30]   (0,0) rectangle (1,1); \node at (0.5,0.5){1};
      \fill[rgreen!60] (1,0) rectangle (2,1); \node at (1.5,0.5){12};
    \end{tikzpicture}
  \end{columns}
\end{columns}
\end{frame}

%=========================== SLIDE 8 - SMOTE + GAUSSIEN =======================
\kick{Génération de données synthétiques}
\begin{frame}{SMOTE + Bruit Gaussien}
\timing{2 min}
\begin{columns}[T]
\column{0.54\textwidth}
  \begin{block}{\textcolor{navydark}{1.\ SMOTE -- Synthetic Minority Over-sampling}}
    \small Interpolation entre \textbf{voisins proches (k-NN)} : un point est généré sur le segment vers un voisin.
    \[ x_{new} = x_i + \lambda\,(x_{voisin}-x_i),\quad \lambda\in[0,1] \]
  \end{block}
  \vspace{0.3em}
  \begin{block}{\textcolor{navydark}{2.\ Ajout de bruit gaussien}}
    \small Un léger bruit normal évite des données trop « lisses » et limite le sur-apprentissage.
    \[ x_{final} = x_{smote} + \varepsilon,\quad \varepsilon\sim\mathcal{N}(0,\sigma^2) \]
    {\footnotesize\color{muted}$\sigma = 0.01 \times$ écart-type de la feature.}
  \end{block}
\column{0.46\textwidth}
  \centering {\color{gold}\small Interpolation SMOTE}\\[2pt]
  \begin{tikzpicture}[scale=0.5]
    \draw[gold,dashed] (1,1)--(3,5); \draw[gold,dashed] (3,5)--(7,4.5);
    \foreach \p in {(1,1),(3,5),(5,3),(7,4.5),(4,0.5)}
      \fill[rred] \p circle (5pt);
    \foreach \p in {(1.8,2.6),(2.4,4.2),(5,4.75)}
      \fill[gold] \p circle (4pt);
  \end{tikzpicture}\\[4pt]
  {\color{gold}\small Bruit gaussien $\mathcal{N}(0,\sigma^2)$}\\[-4pt]
  \begin{tikzpicture}
    \begin{axis}[rgaxis,height=2.5cm,ticks=none,domain=-3:3,samples=60,
      axis lines=none,ymin=0]
      \addplot[gold,line width=1.2pt]{exp(-x^2/0.7)};
    \end{axis}
  \end{tikzpicture}
\end{columns}
\end{frame}

%=========================== SLIDE 9 - 4 CLASSES ==============================
\kick{Multi-classification}
\begin{frame}{4 classes de risque}
\timing{30 s}
\vspace{0.2em}
{\color{muted} Le problème binaire devient une échelle de risque, plus utile qu'un simple oui/non.}
\vspace{1em}

\newcommand{\riskcard}[5]{%
  \begin{tikzpicture}
    \node[fill=#1,rounded corners=8pt,inner sep=8pt,text width=0.9\linewidth,
          text=#5,minimum height=2.6cm,align=left]{%
      {\bfseries\small #3}\\[4pt]{\Large\bfseries #2}\\[4pt]{\small #4}};
  \end{tikzpicture}}
\begin{columns}[T]
  \column{0.25\textwidth}\riskcard{rgreen}{Faible}{0--25\%}{Entreprise saine}{navydark}
  \column{0.25\textwidth}\riskcard{ryellow}{Modéré}{25--50\%}{Vigilance requise}{navydark}
  \column{0.25\textwidth}\riskcard{rorange}{Élevé}{50--75\%}{Surveillance active}{navydark}
  \column{0.25\textwidth}\riskcard{rred}{Critique}{75--100\%}{Risque imminent}{white}
\end{columns}

\vspace{1.2em}
\begin{tikzpicture}
  \shade[left color=rgreen,right color=rred,middle color=rorange]
    (0,0) rectangle (\linewidth,0.5);
\end{tikzpicture}
\end{frame}

%=========================== SLIDE 10 - BENCHMARK =============================
\kick{Comparaison des modèles ML}
\begin{frame}{Benchmark des modèles testés}
\timing{2 min}
\begin{columns}[T]
\column{0.56\textwidth}
  \footnotesize
  \renewcommand{\arraystretch}{1.25}
  \arrayrulecolor{navy}
  \begin{tabular}{lcccc}
    \toprule
    \rowcolor{navy}
    \textcolor{gold}{\bfseries Modèle} & \textcolor{gold}{Acc.} & \textcolor{gold}{F1} & \textcolor{gold}{AUC} & \textcolor{gold}{Train}\\
    \midrule
    Régression Log. & 78\% & 0.72 & 0.81 & Très rapide\\
    Random Forest   & 89\% & 0.86 & 0.92 & Moyen\\
    SVM (RBF)       & 82\% & 0.78 & 0.85 & Lent\\
    KNN             & 80\% & 0.75 & 0.83 & Rapide\\
    \rowcolor{gold!18}
    \textbf{XGBoost} $\star$ & \textbf{94\%} & \textbf{0.93} & \textbf{0.97} & Moyen\\
    LightGBM        & 93\% & 0.91 & 0.96 & Rapide\\
    \bottomrule
  \end{tabular}\\[0.8em]
  {\color{gold}\small $\star$ XGBoost domine sur toutes les métriques.}
\column{0.44\textwidth}
  \centering {\color{gold}\small F1-Score par modèle}\\[2pt]
  \begin{tikzpicture}
    \begin{axis}[rgaxis,height=4.6cm,xbar,xmin=0,xmax=1,
      symbolic y coords={XGBoost,LightGBM,RF,SVM,KNN,RegLog},
      ytick=data,xtick=\empty,
      nodes near coords,nodes near coords style={text=ink,font=\scriptsize}]
      \addplot[fill=rbluel,draw=none] coordinates
        {(0.75,RegLog)(0.78,SVM)(0.75,KNN)(0.86,RF)(0.91,LightGBM)};
      \addplot[fill=gold,draw=none] coordinates {(0.93,XGBoost)};
    \end{axis}
  \end{tikzpicture}
\end{columns}
\end{frame}

%=========================== SLIDE 11 - POURQUOI XGBOOST ======================
\kick{Justification du choix}
\begin{frame}{Pourquoi XGBoost ?}
\timing{1 min 30}
\begin{columns}[T]
\column{0.55\textwidth}
  \small
  \begin{enumerate}
    \item \textbf{Performance} : meilleur F1 \& AUC-ROC.
    \item \textbf{Déséquilibre} géré : \texttt{scale\_pos\_weight}.
    \item \textbf{Régularisation L1+L2} : anti-surapprentissage.
    \item \textbf{Robuste aux outliers} (arbres de décision).
    \item \textbf{Feature importance} $\rightarrow$ interprétabilité.
    \item \textbf{Optuna} : tuning bayésien (30 essais).
    \item \textbf{Inférence} $<$ 100 ms par requête.
    \item \textbf{Écosystème mature}, standard en finance.
  \end{enumerate}
\column{0.45\textwidth}
  \centering {\color{gold}\small Gradient Boosting séquentiel}\\[4pt]
  \begin{tikzpicture}[scale=0.8,node distance=1.1cm,
    tree/.style={circle,fill=rbluel,text=navydark,font=\footnotesize\bfseries,minimum size=0.8cm}]
    \node[tree](t1){T1};
    \node[tree,right=of t1](t2){T2};
    \node[tree,right=of t2](t3){T3};
    \node[tree,fill=gold,right=of t3](t4){T4};
    \draw[gold,-{Latex}](t1)--(t2);
    \draw[gold,-{Latex}](t2)--(t3);
    \draw[gold,-{Latex}](t3)--(t4);
  \end{tikzpicture}\\[6pt]
  \rgcard{0.92\linewidth}{%
    {\color{gold}\faTimesCircle\ \bfseries Pourquoi pas les autres ?}\\[3pt]
    \scriptsize
    \textbullet\ Rég. logistique : pas de non-linéarités.\\
    \textbullet\ SVM : lent, peu interprétable.\\
    \textbullet\ Random Forest : sans boosting séquentiel.\\
    \textbullet\ Deep Learning : trop complexe, « boîte noire ».}
\end{columns}
\end{frame}

%=========================== SLIDE 12 - OPTUNA ================================
\kick{Optimisation des hyperparamètres}
\begin{frame}{Tuning avec Optuna}
\timing{1 min}
\begin{columns}[T]
\column{0.52\textwidth}
  \begin{itemize}
    \item \textbf{Optimisation bayésienne} -- plus intelligente que Grid Search.
    \item \textbf{30 essais} d'optimisation automatique.
    \item \textbf{Validation croisée stratifiée (5-fold)}.
  \end{itemize}
  \vspace{0.4em}
  \rgcard{0.93\linewidth}{%
    {\color{gold}\bfseries Hyperparamètres optimisés}\\[4pt]
    \small \texttt{max\_depth} \textbullet\ \texttt{learning\_rate} \textbullet\ \texttt{n\_estimators}
    \textbullet\ \texttt{subsample} \textbullet\ \texttt{colsample\_bytree}
    \textbullet\ \texttt{reg\_alpha} \textbullet\ \texttt{reg\_lambda}}
\column{0.48\textwidth}
  \centering {\color{gold}\small Convergence des essais}\\[2pt]
  \begin{tikzpicture}
    \begin{axis}[rgaxis,height=4.6cm,xlabel={essais},ylabel={AUC-ROC},
      xmin=0,xmax=30,ymin=0.6,ymax=1,ymajorgrids]
      \addplot[only marks,mark=*,mark size=1pt,rbluel] coordinates
        {(1,.70)(2,.66)(4,.74)(6,.76)(8,.79)(10,.80)(12,.85)(14,.87)(16,.88)
         (18,.90)(20,.905)(22,.91)(24,.918)(26,.92)(28,.925)(30,.929)};
      \addplot[gold,line width=1.4pt] coordinates
        {(1,.70)(4,.74)(6,.78)(10,.84)(14,.88)(18,.905)(22,.915)(26,.925)(30,.93)};
    \end{axis}
  \end{tikzpicture}
\end{columns}
\end{frame}

%=========================== SLIDE 13 - ARCHITECTURE =========================
\kick{Architecture technique}
\begin{frame}{Architecture de la plateforme}
\timing{1 min 30}
\vspace{0.3em}
\centering
\begin{tikzpicture}[every node/.style={font=\small},
  tier/.style={fill=card,draw=gold,line width=1pt,rounded corners=8pt,
               text width=0.86\linewidth,inner sep=8pt,text=ink},
  arr/.style={gold,-{Latex},line width=1.2pt}]
  \node[tier](fe){%
    {\color{gold}\faDesktop\ \bfseries Frontend} \quad{\color{muted}\footnotesize HTML / CSS / JS pur}\\[2pt]
    \footnotesize Interface responsive \textbullet\ Dashboard KPIs \& graphiques \textbullet\ Historique \textbullet\ Mode démo (15 entreprises MA)};
  \node[tier,below=0.9cm of fe](be){%
    {\color{gold}\faCogs\ \bfseries Backend} \quad{\color{muted}\footnotesize FastAPI -- Python}\\[2pt]
    \footnotesize API RESTful (8 endpoints) \textbullet\ Validation Pydantic \textbullet\ Rapports PDF \textbullet\ SQLite (historique)};
  \node[tier,below=0.9cm of be](ml){%
    {\color{gold}\faBrain\ \bfseries ML Pipeline} \quad{\color{muted}\footnotesize XGBoost -- joblib}\\[2pt]
    \footnotesize Préprocessing $\rightarrow$ Feature selection $\rightarrow$ XGBoost $\rightarrow$ SHAP \textbullet\ Prédiction multi-classes + proba};
  \draw[arr](fe)--node[right,muted,font=\scriptsize]{REST / JSON}(be);
  \draw[arr](be)--node[right,muted,font=\scriptsize]{appel modèle}(ml);
\end{tikzpicture}
\end{frame}

%=========================== SLIDE 14 - FONCTIONNALITES ======================
\kick{Fonctionnalités clés}
\begin{frame}{Fonctionnalités}
\timing{1 min}
\vspace{0.3em}
\begin{columns}[T]
\column{0.33\textwidth}
  \feat{\faSearch}{Analyse unitaire}{Saisie manuelle des ratios.}\\[6pt]
  \feat{\faChartLine}{Dashboard}{Distribution \& tendances.}
\column{0.33\textwidth}
  \feat{\faTable}{Analyse batch}{Upload CSV / Excel.}\\[6pt]
  \feat{\faHistory}{Historique}{Filtrable et paginé.}
\column{0.33\textwidth}
  \feat{\faFilePdf}{Export PDF}{Rapport + recommandations.}\\[6pt]
  \feat{\faBullseye}{Interprétabilité}{SHAP values par prédiction.}
\end{columns}
\vspace{1em}
\centering
\tikz\node[fill=gold!15,draw=golddeep,rounded corners=8pt,inner sep=8pt,text=goldsoft,font=\small\bfseries]
  {\faBolt\ Temps réel -- réponse en moins de 100 ms par requête.};
\end{frame}

%=========================== SLIDE 15 - DEMO =================================
\kick{Démonstration en direct}
\begin{frame}{Place à la démo}
\centering
\vspace{0.4em}
{\color{gold}\Huge\faPlayCircle}\\[1em]
\begin{columns}[T]
\column{0.45\textwidth}
  \feat{\faCircle[regular]}{Scénario 1}{Entreprise à risque faible.}\\[6pt]
  \feat{\faCircle[regular]}{Scénario 3}{Upload batch d'un fichier CSV.}
\column{0.45\textwidth}
  \feat{\faCircle[regular]}{Scénario 2}{Entreprise à risque critique.}\\[6pt]
  \feat{\faCircle[regular]}{Rapport PDF}{Génération \& export.}
\end{columns}
\vspace{1em}
\tikz\node[fill=gold!12,draw=golddeep,rounded corners=6pt,inner sep=8pt,text width=0.85\linewidth,text=ink,font=\small]
  {{\color{gold}\bfseries Note présentateur :} préparer le backend lancé
   (\texttt{uvicorn main:app --port 8000}) + le fichier \texttt{test\_batch.csv}.};
\end{frame}

%=========================== SLIDE 16 - RESULTATS ============================
\kick{Résultats du modèle final}
\begin{frame}{Résultats \& Métriques}
\timing{1 min}
\centering
\newcommand{\statcard}[2]{%
  \begin{tikzpicture}
    \node[fill=gold!12,draw=golddeep,rounded corners=8pt,inner sep=8pt,
          text width=0.85\linewidth,align=center]{%
      {\color{gold}\Large\bfseries #1}\\[2pt]{\color{muted}\footnotesize #2}};
  \end{tikzpicture}}
\begin{columns}[c]
  \column{0.3\textwidth}\statcard{94\%}{Accuracy}
  \column{0.3\textwidth}\statcard{0.93}{F1-Score macro}
  \column{0.3\textwidth}\statcard{0.97}{AUC-ROC}
\end{columns}
\vspace{0.8em}
\begin{columns}[T]
\column{0.32\textwidth}\centering
  {\color{gold}\scriptsize Matrice de confusion}\\[2pt]
  \begin{tikzpicture}[scale=0.62,every node/.style={font=\scriptsize\bfseries,text=white}]
    \fill[rgreen!75] (0,1) rectangle (1,2); \node at (0.5,1.5){1577};
    \fill[rred!30]   (1,1) rectangle (2,2); \node at (1.5,1.5){43};
    \fill[rred!30]   (0,0) rectangle (1,1); \node at (0.5,0.5){33};
    \fill[rgreen!65] (1,0) rectangle (2,1); \node at (1.5,0.5){311};
  \end{tikzpicture}
\column{0.30\textwidth}\centering
  {\color{gold}\scriptsize Courbe ROC}\\[2pt]
  \begin{tikzpicture}[scale=0.62]
    \draw[muted] (0,0) rectangle (2.6,2.6);
    \draw[muted,dashed] (0,0)--(2.6,2.6);
    \draw[gold,line width=1.4pt]
      (0,0) .. controls (0.1,1.6) and (0.4,2.4) .. (1,2.5) -- (2.6,2.6);
    \node[gold,font=\scriptsize\bfseries] at (1.7,0.6){AUC 0.97};
  \end{tikzpicture}
\column{0.36\textwidth}\centering
  {\color{gold}\scriptsize Top 5 features}\\[2pt]
  \begin{tikzpicture}
    \begin{axis}[rgaxis,height=3.0cm,xbar,xmin=0,xmax=0.45,
      symbolic y coords={Retained,WorkCap,Current,ROA,Debt},
      ytick=data,xtick=\empty,yticklabel style={font=\tiny}]
      \addplot[fill=rbluel,draw=none] coordinates
        {(0.10,Retained)(0.12,WorkCap)(0.19,Current)(0.30,ROA)};
      \addplot[fill=gold,draw=none] coordinates {(0.39,Debt)};
    \end{axis}
  \end{tikzpicture}
\end{columns}
\end{frame}

%=========================== SLIDE 17 - CONCLUSION ===========================
\kick{Bilan}
\begin{frame}{Conclusion}
\timing{30 s}
\vspace{0.6em}
\large
\begin{itemize}
  \setbeamertemplate{itemize item}{\color{rgreen}\faCheckCircle}
  \item Pipeline ML \textbf{complet de bout en bout}.
  \item Résolution du \textbf{déséquilibre} via SMOTE + bruit gaussien.
  \item \textbf{XGBoost} : choix justifié par performances et interprétabilité.
  \item Plateforme web \textbf{fonctionnelle et intuitive}.
  \item Applicable au \textbf{contexte bancaire marocain}.
\end{itemize}
\end{frame}

%=========================== SLIDE 18 - PERSPECTIVES =========================
\kick{Et après ?}
\begin{frame}{Perspectives \& Améliorations}
\timing{30 s}
\vspace{0.3em}
\begin{columns}[T]
\column{0.5\textwidth}
  \feat{\faUniversity}{Données réelles marocaines}{Bank Al-Maghrib, OMPIC.}\\[6pt]
  \feat{\faFileAlt}{Données textuelles (NLP)}{Rapports d'audit analysés.}\\[6pt]
  \feat{\faCloud}{Déploiement cloud}{AWS / Azure en production.}
\column{0.5\textwidth}
  \feat{\faChartArea}{Monitoring du modèle}{Détection de drift.}\\[6pt]
  \feat{\faBuilding}{Extension PME / TPE}{Features adaptées.}\\[6pt]
  \feat{\faPlug}{API publique}{Intégration tierce (SaaS).}
\end{columns}
\end{frame}

%=========================== SLIDE 19 - MERCI ================================
\kick{}
{
\setbeamertemplate{footline}{}
\begin{frame}[plain]
\begin{tikzpicture}[remember picture,overlay]
  \fill[navydark](current page.south west) rectangle (current page.north east);
\end{tikzpicture}
\centering
\vspace{1.4cm}
{\color{gold}\Huge\faShieldAlt}\\[0.8em]
{\color{white}\huge\bfseries Merci pour votre attention}\\[0.4em]
{\color{gold}\Large\bfseries Questions ?}\\[1.2em]
\begin{tabular}{ccc}
  {\color{white}\bfseries Étudiant(s)} & {\color{white}\bfseries Contact} & {\color{white}\bfseries GitHub}\\[2pt]
  {\color{goldsoft}\itshape [ à compléter ]} & {\color{goldsoft}\itshape [ email ]} & {\color{goldsoft}achrafbaaiz/riskguard-pro}\\
\end{tabular}\\[1.2em]
\colorbox{white}{\qrcode[height=2.2cm]{https://github.com/achrafbaaiz/riskguard-pro}}\\[4pt]
{\color{muted}\footnotesize QR $\rightarrow$ dépôt GitHub du projet}
\end{frame}
}

\end{document}
